%% jss-lint: Automated Style Checking for JSS Manuscripts
%% Manuscript for submission to the Journal of Statistical Software
%%
%% NOTE: Compilation requires jss.cls, downloadable from
%%       https://www.jstatsoft.org/style

\documentclass[article]{jss}
\usepackage{longtable}

%% -- Article metainformation -----------------------------------------

\author{Manuel Koller\\Author Institution}
\Plainauthor{Manuel Koller}

\title{\pkg{jss-lint}: Automated Style Checking for
  \emph{Journal of Statistical Software} Manuscripts}
\Plaintitle{jss-lint: Automated Style Checking for
  Journal of Statistical Software Manuscripts}

\Abstract{%
  Authors submitting manuscripts to the \emph{Journal of Statistical
  Software} (JSS) must satisfy a comprehensive style guide covering
  preamble macros, semantic markup, citation conventions,
  capitalisation, code formatting, bibliography completeness, and
  typography. Style violations that slip through to peer review create
  unnecessary friction and additional revision rounds. We present
  \pkg{jss-style-checker}, an open-source \proglang{Python}~3.10+
  command-line tool and library providing deterministic, rule-based
  validation of \LaTeX{} manuscripts against the JSS style guide. The
  tool parses manuscripts using an abstract syntax tree (AST) approach
  via \pkg{pylatexenc}, enabling semantic analysis that avoids false
  positives from comment-only or verbatim occurrences. Fifty-three
  rules across 16 categories are implemented, covering \code{.tex},
  \code{.Rnw} (Sweave/Knitr), and \code{.Rmd} (\proglang{R}~Markdown)
  source formats. A \code{--fix} flag applies deterministic
  auto-corrections in place; a \code{--dry-run} flag previews them.
  A journal plugin system based on \proglang{Python} entry points
  enables other journals to build on the core infrastructure without
  modifying the package.

  The tool was validated against a corpus of 493 real JSS-format
  manuscripts drawn from CRAN package vignettes. After sixteen
  iterations of a precision evaluation-improvement loop, all thirty
  rules with at least ten reviewed violations achieve at least 90\%
  precision (overall: 98.2\% across 6{,}739 reviewed instances). The
  development methodology and retrospective lessons are described in
  detail and offer a replicable template for similar tools. The
  package is available from the \proglang{Python} Package Index as
  \pkg{jss-style-checker}; the present manuscript was validated with
  the tool before submission.%
}

\Keywords{%
  \LaTeX{}, style checking, academic writing, linting,
  \proglang{Python}, \emph{Journal of Statistical Software},
  large language models%
}
\Plainkeywords{%
  LaTeX, style checking, academic writing, linting, Python,
  Journal of Statistical Software, AI-assisted development,
  large language models%
}

\Address{
  Manuel Koller\\
  Department\\
  Institution\\
  E-mail: \email{author@institution.edu}
}

\begin{document}

%%--------------------------------------------------------------------
\section{Introduction}
\label{sec:intro}
%%--------------------------------------------------------------------

The \emph{Journal of Statistical Software} (JSS) publishes software
packages and statistical methods alongside reproducible examples, and
maintains a detailed style guide governing manuscript preparation
\citep{jss_styleguide}. The guide addresses a wide range of
requirements covering \LaTeX{} preamble macros, semantic markup
conventions for programming language and package names, citation style
(natbib), capitalisation in section headings and bibliographic entries,
line-width limits in code listings, label requirements for floating
environments, and typography details such as punctuation after Latin
abbreviations.

While each individual requirement is straightforward, satisfying all
of them simultaneously across a long manuscript is challenging,
especially for authors who are unfamiliar with JSS conventions or
whose workflow evolved from another journal. Checking compliance
manually during revision is time-consuming. Automated tools exist for
general \LaTeX{} checking (\pkg{chktex}~\citep{chktex},
\pkg{lacheck}~\citep{lacheck}, \pkg{textlint}~\citep{texlint_js}),
but none implement JSS-specific rules.

We present \pkg{jss-style-checker} (command-line entry point
\code{jss-lint}), which to our knowledge is the first dedicated
automated style checker for JSS manuscripts. The tool takes one or
more \code{.tex}, \code{.Rnw}, \code{.Rmd}, and \code{.bib} files as
input and produces a structured violation report. Its design emphasises
three properties:

\begin{enumerate}
  \item \emph{Determinism}: identical input always produces identical
    output; there are no probabilistic or machine-learning components.
  \item \emph{Semantic analysis}: the tool parses \LaTeX{} into an
    AST rather than applying regular expressions to raw text, avoiding
    false positives inside comments or verbatim blocks.
  \item \emph{Extensibility}: a journal plugin system based on
    \proglang{Python} entry points allows other journals to define
    custom rule sets without modifying the core package.
\end{enumerate}

Empirical validation on 493 real JSS-format manuscripts confirms that
the approach is effective: all thirty rules with sufficient reviewed
data achieve at least 90\% precision, and the overall precision across
6{,}739 reviewed violation instances is 98.2\%. This paper also
documents the development process in detail, including the iterative
evaluation-improvement methodology that drove rule precision from
initial implementations to the published release.

This paper is organised as follows. Section~\ref{sec:styleguide}
surveys the JSS style requirements. Section~\ref{sec:design} describes
the design and architecture. Section~\ref{sec:aidev} documents the
AI-assisted development methodology, including the iterative
evaluation-improvement loop. Section~\ref{sec:lessons} reflects
candidly on what worked, what proved difficult, and what we would do
differently. Section~\ref{sec:rules} documents each implemented rule
category. Section~\ref{sec:eval} presents empirical precision results
on 493 real manuscripts. Section~\ref{sec:usage} covers installation
and usage. Section~\ref{sec:extensibility} explains how to create
journal plugins. Section~\ref{sec:comparison} compares
\pkg{jss-style-checker} with related tools, and
Section~\ref{sec:conclusions} concludes.

%%--------------------------------------------------------------------
\section{JSS style guide}
\label{sec:styleguide}
%%--------------------------------------------------------------------

The JSS style guide \citep{jss_styleguide} is the authoritative
reference for manuscript preparation. Authors are expected to study
the guide and the companion template before submitting. This section
surveys the requirements that \pkg{jss-style-checker} enforces and
explains why manual compliance is difficult.

\subsection{Overview of requirements}

The JSS \LaTeX{} document class (\code{jss.cls}) defines seven
semantic markup commands that authors must use consistently:

\begin{itemize}
  \item \code{\textbackslash proglang\{name\}}: programming languages
    (e.g.,{} \proglang{R}, \proglang{Python}, \proglang{Julia}).
  \item \code{\textbackslash pkg\{name\}}: software packages
    (e.g.,{} \pkg{ggplot2}, \pkg{numpy}, \pkg{lme4}).
  \item \code{\textbackslash code\{snippet\}}: inline code,
    filenames, function names, and command-line options.
  \item \code{\textbackslash email\{address\}}: e-mail addresses.
  \item \code{\textbackslash file\{path\}}: file paths.
  \item \code{\textbackslash fct\{name\}}: function names (optional
    alternative to \code{\textbackslash code}).
  \item \code{\textbackslash url\{uri\}}: hyperlinks (with
    \code{\textbackslash href} for labeled links).
\end{itemize}

Beyond markup commands, the guide specifies six required preamble
macros (\code{\textbackslash Plainauthor},
\code{\textbackslash Plaintitle},
\code{\textbackslash Abstract},
\code{\textbackslash Keywords},
\code{\textbackslash Plainkeywords},
\code{\textbackslash Address}) that provide plain-text metadata for
the journal management system. All six must appear before
\code{\textbackslash begin\{document\}}.

Citation style follows natbib \citep{natbib}: parenthetical references
use \code{\textbackslash citep\{key\}}, and textual references use
\code{\textbackslash citet\{key\}}. The style guide specifically
prohibits brackets-within-brackets constructs such as
\code{(\textbackslash cite\{key\})}, which produce
doubly-parenthesised output; natbib commands should be used
throughout.

Section headings use sentence case (only the first word and proper
nouns capitalised), which differs from the title case common in other
publication venues. BibTeX title fields, conversely, must be in title
case and must protect capitalisation with curly braces
(e.g.,{} \code{title = \{An Overview of \{R\}\}}) to prevent BibTeX
styles from down-casing proper nouns.

Code examples inside verbatim environments
(\code{verbatim}, \code{Code}, \code{CodeInput}, \code{CodeOutput},
\code{Sinput}, \code{Soutput}, \code{lstlisting}) must fit within
the column width of the typeset page. All floating environments
(figures, tables, equations) must carry a
\code{\textbackslash label\{\}} for cross-referencing. Typography
rules require trailing commas: the forms ``e.g.,'' and ``i.e.,'' are
mandated (not the bare forms), and the \LaTeX{} macro
\code{\textbackslash top} is required for the matrix transpose symbol
rather than the common but ambiguous \verb|^T|.

\subsection{The case for automated checking}

We argue that style compliance checking and scientific peer review are
better kept separate. Peer review is intended to assess the scientific
contribution, evaluate methodology, and identify gaps in the
literature; style conformance is a separate, mechanical concern with
deterministic remedies. Violations that can be detected and fixed
before submission need not consume any reviewer attention. Automated
pre-submission checking addresses this class of issues completely, for
the rules it covers, before a manuscript enters the review pipeline.

\subsection{Motivating example}
\label{sec:motivating-example}

Consider the following non-compliant \LaTeX{} fragment:

\begin{verbatim}
% Non-compliant example (do not copy)
\section{R Statistical Computing}

We fit a linear model in Python using the
statsmodels package (\cite{rcore2024}).
\end{verbatim}

This fragment violates three JSS rules simultaneously: the section
heading uses title case instead of sentence case
(\code{JSS-CAP-002}); the language name \proglang{Python} and the
package name \pkg{statsmodels} appear without semantic markup
(\code{JSS-MARKUP-001}/\code{JSS-MARKUP-002}); and a bare
\code{\textbackslash cite\{\}} inside parentheses creates a
brackets-within-brackets construct (\code{JSS-CITE-001}). Running
\code{jss-lint} on a file containing this fragment produces:

\begin{verbatim}
example.tex
  2:1  warning  [JSS-CAP-002]  Title Case in section heading
     -> Use sentence case: "R statistical computing"
  5:29  warning  [JSS-MARKUP-001]  Bare language name "Python"
     -> Replace with \proglang{Python}
  5:45  warning  [JSS-MARKUP-002]  Bare package "statsmodels"
     -> Replace with \pkg{statsmodels}
  6:28  error    [JSS-CITE-001]  Bare \cite{} inside parentheses
     -> Use \citep{rcore2024} or \citet{rcore2024}

4 violations (1 error, 3 warnings)
\end{verbatim}

The corrected fragment is:

\begin{Code}
\section{R statistical computing}

We fit a linear model in \proglang{Python} using the
\pkg{statsmodels} package \citep{rcore2024}.
\end{Code}

%%--------------------------------------------------------------------
\section{Design and architecture}
\label{sec:design}
%%--------------------------------------------------------------------

\subsection{Design philosophy}
\label{sec:philosophy}

\pkg{jss-style-checker} is built around three principles.

\textbf{AST-based analysis.} The tool parses \LaTeX{} source files
into an abstract syntax tree using \pkg{pylatexenc}
\citep{pylatexenc} rather than applying regular expressions to raw
text. AST-based analysis correctly identifies the semantic role of
each token: a bare word \verb|Python| inside a \code{verbatim}
environment is not a markup violation because verbatim text is
rendered as-is; a bare \verb|Python| in a paragraph body is a
violation. Making this distinction with regular expressions requires
tracking nested environments and comment state explicitly,
substantially increasing implementation complexity and the risk of
false positives.

\textbf{Non-fatal parse errors.} Unusual or non-standard \LaTeX{}
that \pkg{pylatexenc} cannot fully parse does not crash the tool.
Parse warnings are stored alongside any violations found so that
partial checking is always possible, even on unconventional
manuscripts.

\textbf{Single-file scope.} \pkg{jss-style-checker} analyses each
input file independently. \code{\textbackslash input\{\}} and
\code{\textbackslash include\{\}} directives are \emph{not} followed;
referenced sub-files are ignored unless they are passed explicitly on
the command line. Authors with multi-file manuscripts should pass all
component files (and the shared \code{.bib} file) as arguments in a
single invocation:

\begin{CodeInput}
jss-lint main.tex appendix.tex refs.bib
\end{CodeInput}

\textbf{Deterministic output.} Given the same input file, the tool
always produces the same set of violations in the same order (sorted
by file, line, and column). This property is essential for integrating
the tool into continuous integration (CI) pipelines or pre-commit
hooks.

\subsection{Package architecture}
\label{sec:architecture}

\pkg{jss-style-checker} follows a \code{src/} layout with a clear
separation between the public API, the rule engine, journal
implementations, and output formatters. The top-level structure is:

\begin{Code}
src/texlint/
    api.py         # Public types
    cli.py         # Click CLI entry point
    config.py      # Configuration loader
    core/
        parser.py  # source-file parsers (tex, bib, rnw, rmd)
        engine.py  # load_journal, run
    journals/
        jss/       # JSS rule module
    output/        # terminal, json, html formatters
\end{Code}

The core pipeline is:

\begin{enumerate}
  \item \textbf{Parse}: \code{parse\_tex\_file(path)},
    \code{parse\_rnw\_file(path)}, and
    \code{parse\_rmd\_file(path)} produce \code{ParsedTexFile} and
    \code{ParsedRmdFile} objects respectively;
    \code{parse\_bib\_file(path)} produces a \code{ParsedBibFile}.
    Each carries parsed content and any parse-time violations.
  \item \textbf{Assemble}: the parsed files are collected into a
    \code{ParsedDocument} container.
  \item \textbf{Run}: \code{run(doc, config)} loads the journal
    module, executes all rules against the document, filters
    suppressed rule IDs, and returns a \code{ComplianceReport}.
  \item \textbf{Format}: the report is rendered by the appropriate
    formatter (terminal, JSON, or HTML) and written to standard
    output.
\end{enumerate}

The principal public types are summarised in
Table~\ref{tab:types}.

\begin{table}[ht]
  \centering
  \begin{tabular}{lp{9cm}}
    \hline
    Type & Description \\
    \hline
    \code{Violation} &
      Frozen dataclass: file, line, column, rule ID, severity
      (\code{"error"} or \code{"warning"}), message,
      optional human-readable suggestion, and an optional
      \code{fix} (\code{FixSuggestion}) carrying a character-offset
      patch for auto-correction. \\
    \code{Rule} &
      Frozen dataclass: rule ID, description, optional
      \code{formats} set, and a check function
      \code{(ParsedDocument) -> Iterator[Violation]}. \\
    \code{RuleCategory} &
      Named grouping of one or more \code{Rule} objects. \\
    \code{ComplianceReport} &
      Top-level result: journal, mode, version, list of
      violations, per-category summaries, compliance percentage. \\
    \code{ToolConfig} &
      Immutable merged configuration: journal ID, mode,
      output format, frozen set of ignored rule IDs. \\
    \code{ParsedDocument} &
      Container for all \code{ParsedTexFile},
      \code{ParsedBibFile}, and \code{ParsedRmdFile} objects. \\
    \code{ParsedRmdFile} &
      Parsed \code{.Rmd} file: YAML frontmatter dict, list of
      \code{RmdHeading}, \code{RmdProseBlock}, and
      \code{RmdCodeBlock} objects, inline
      \code{ParsedTexFile} fragments. \\
    \hline
  \end{tabular}
  \caption{Principal public types in the \pkg{jss-style-checker}
    API (\code{src/texlint/api.py}).}
  \label{tab:types}
\end{table}

\subsection{Configuration system}
\label{sec:config}

Configuration follows a three-tier precedence model: built-in defaults
are overridden by a \code{.jss-lint.toml} file in the working
directory, which is in turn overridden by CLI flags. The TOML
configuration file recognises four keys:

\begin{Code}
journal = "jss"       # rule-set to apply
mode = "author"       # author | reviewer
output = "terminal"   # terminal | json | html
ignore-rules = [      # suppress specific rule IDs
    "JSS-CITE-001"
]
\end{Code}

\subsection{Output modes}
\label{sec:output-modes}

Two audience modes are provided.

\textbf{Author mode} produces a flat, annotated list of violations
sorted by location. Each violation displays its line and column,
severity, rule ID, a human-readable message, and a concrete suggestion
for how to fix it. This mode is intended for use during manuscript
preparation.

\textbf{Reviewer mode} produces a compliance summary table showing
pass/fail status and violation count for each rule category, together
with an overall compliance percentage. This mode is intended for
editorial workflows.

Three output formats are supported:

\begin{itemize}
  \item Terminal: coloured output using \pkg{rich}~\citep{rich},
    suitable for interactive use.
  \item JSON: machine-readable output with a versioned schema,
    suitable for CI/CD integration or programmatic post-processing.
  \item HTML: self-contained report generated with
    \pkg{jinja2}~\citep{jinja2} templates, suitable for sharing with
    collaborators.
\end{itemize}

The \code{--verbose} flag (\code{-v}) adds a list of rules that were
skipped because they do not apply to the input format (e.g.,{} \LaTeX{}
preamble rules when only \code{.Rmd} files are provided). In JSON
output the skipped rules are included as a \code{skipped\_rules} array.

\subsection{R source file support}
\label{sec:r-source}

Many JSS authors write their manuscripts in \proglang{R} source formats
that mix prose and code before compilation to \LaTeX{}.
\pkg{jss-style-checker} accepts these formats directly.

\textbf{Sweave/Knitr (\code{.Rnw}).} The parser strips \proglang{R}
code chunks (delimited by \verb|<<...>>=| \ldots{} \verb|@|) by
replacing each chunk line with a blank line, preserving the original
line numbering so that violation line numbers in the \code{.Rnw} file
remain accurate. Inline \verb|\Sexpr{...}| expressions are replaced
with an inert placeholder. The remaining \LaTeX{} text is then parsed
and checked by all rules applicable to \code{.tex} files.

\textbf{\proglang{R}~Markdown (\code{.Rmd}).} The parser separates the
document into three element types: the YAML frontmatter block, fenced
code blocks (\verb|```{r}| \ldots{} \verb|```|), and prose blocks.
Inline \LaTeX{} macros in prose blocks (e.g.,{} \verb|\pkg{ggplot2}|,
\verb|\code{lm()}|) are extracted and parsed with \pkg{pylatexenc}.
Three \code{.Rmd}-specific rules are provided
(\code{JSS-RMD-PREAMBLE-001}, \code{JSS-RMD-CAP-002},
\code{JSS-RMD-CITE-002}), and several existing rules
(\code{JSS-WIDTH-001}, \code{JSS-CODE-002}, \code{JSS-CODE-003},
\code{JSS-ABBR-001}, \code{JSS-STYLE-001}, \code{JSS-RCODE-001})
are extended to operate on \code{.Rmd} code and prose blocks.

Rules are tagged with a \code{formats} field
(\code{"tex"}/\code{"rnw"} or \code{"rmd"}). When the engine
encounters a rule whose \code{formats} set has no intersection with
the active input formats, that rule is recorded as skipped (reported
with \code{--verbose}).


%%--------------------------------------------------------------------
\section{AI-assisted development methodology}
\label{sec:aidev}
%%--------------------------------------------------------------------

This section documents the development methodology in detail, both for
transparency about how the tool was built and as a template for others
constructing similar rule-based validation tools. The development
combined three distinct practices: a specification-driven workflow,
test-driven rule implementation with coverage requirements, and an
iterative precision-improvement loop grounded in real-world data.
Readers interested primarily in using the tool may skip to
Section~\ref{sec:lessons} or Section~\ref{sec:comparison}.

\subsection{Overview}
\label{sec:aidev-overview}

The tool was developed over approximately two weeks with Claude
Sonnet 4.6 \citep{claude2024} as the primary coding collaborator.
Development produced 37 commits across five feature increments, a
test suite of 611 tests, and sixteen iterations of the precision
evaluation loop. Two broad phases structured the work:

\begin{enumerate}
  \item \emph{Construction phase}: implement rules, tests, and
    infrastructure against synthetic examples, guided by the JSS
    style guide.
  \item \emph{Validation phase}: test against real manuscripts,
    diagnose false positives, and refine rules through targeted fixes.
\end{enumerate}

The construction phase was faster; the validation phase was more
valuable.

\subsection{Specification-driven workflow}
\label{sec:specdriven}

Each feature increment was preceded by a written specification
(\code{spec.md}), an implementation plan (\code{plan.md}), and a
task checklist (\code{tasks.md}), managed with the \pkg{spec-kit}
workflow tool \citep{speckit}. The five increments were:

\begin{enumerate}
  \item Core \LaTeX{} parsing, rule engine, and CLI scaffolding.
  \item Expanded JSS rule set (53 rules across 16 categories).
  \item \proglang{R} source format support (\code{.Rnw} and
    \code{.Rmd}).
  \item Auto-fix infrastructure (\code{--fix}/\code{--dry-run}).
  \item Real-world evaluation loop (\code{eval-jss} CLI and corpus
    tooling).
\end{enumerate}

Writing specifications served two functions that were underappreciated
at the outset. First, they forced early resolution of design decisions
--- such as whether rule severity was per-rule or per-violation ---
before inconsistent implementations could accumulate. Second, they
provided a stable anchor for the AI assistant, reducing the frequency
of responses that drifted from the agreed design. When the AI was
given a clear specification, the quality of generated code was
substantially higher than when it was asked to ``implement the rule''
without additional context.

\subsection{Test-driven rule implementation}
\label{sec:tdd}

Each rule was implemented under a strict test-driven discipline: the
test was written first, confirmed failing, and then the implementation
was written to make it pass. A policy of 100\% branch coverage on all
rule modules was enforced using \pkg{pytest-cov}~\citep{pytest-cov}.
Coverage was measured after each task, and any uncovered branch was
treated as a bug requiring a new test.

This discipline caught several subtle rule interactions that synthetic
positive and negative examples alone would not have revealed. For
example, the branch covering ``content inside a
\code{\textbackslash proglang\{\}} argument'' in
\code{JSS-MARKUP-001} was initially uncovered because no test
included a sentence of the form
\verb|\proglang{Python} has excellent \proglang{R} integration|,
which without the branch would have generated two false positives. The
coverage requirement surfaced this gap immediately.

The final test suite comprises 611 tests: one file per rule category
for unit tests, integration tests for all CLI invocation paths,
\code{.Rnw} and \code{.Rmd} integration tests, and plugin extension
tests using a stub third-party journal module.

\subsection{The precision evaluation-improvement loop}
\label{sec:evalloop}

The most important methodological contribution of this development is
the iterative precision-improvement loop. The loop operates as
follows:

\begin{enumerate}
  \item \textbf{Scan}: run \code{jss-lint} against the full 493-manuscript
    corpus. New violations are inserted into \code{eval.db}; existing
    violations from a previous scan of the same file are superseded.
  \item \textbf{Classify}: route violations to the AI classifier for
    initial TP/FP labelling, then to human review for uncertain or
    high-stakes cases.
  \item \textbf{Report}: compute per-rule precision. Rules below the
    90\% threshold are flagged for investigation.
  \item \textbf{Diagnose}: inspect false-positive violations for each
    flagged rule. Group FPs by root cause --- distinct patterns that
    trigger the rule incorrectly.
  \item \textbf{Fix}: implement a targeted correction for the most
    common root cause. Corrections are deliberately narrow (a single
    additional exclusion condition) to avoid suppressing genuine
    violations.
  \item Repeat from step 1.
\end{enumerate}

Each complete cycle constituted one iteration. The tool reached
publishable precision in sixteen iterations; the first four iterations
established the corpus and infrastructure while scanning incrementally,
and iterations 5--16 focused on diagnosing and fixing rules that fell
below the threshold.

\textbf{Case study: JSS-CODE-003.} The rule \code{JSS-CODE-003},
which flags missing spaces around binary operators or after commas
in code blocks, provides the most instructive example of this loop.
Table~\ref{tab:code003} traces its precision history through
the key iterations.

\begin{table}[ht]
  \centering
  \begin{tabular}{llrrrl}
    \hline
    Iteration & Batch & TP & FP & Precision & Action \\
    \hline
    10 & batch-025 &  72 &  6 & 92.3\% & initial corpus \\
    13 & batch-033 & 775 & 429 & 64.4\% & FP explosion detected \\
    14a & batch-034 & 412 & 223 & 64.9\% & 3 targeted fixes \\
    14b & batch-035 & 485 & 198 & 71.0\% & +2 more fixes \\
    14c & batch-036 & 514 & 191 & 72.9\% & human review queued \\
    15 & batch-037 & 809 &   7 & 99.1\% & 191 AI labels corrected \\
    \hline
  \end{tabular}
  \caption{JSS-CODE-003 precision history across key evaluation
    iterations. The large FP count in iteration~13 arose from scanning
    300 additional manuscripts that exposed patterns not present in
    the initial 200-paper corpus.}
  \label{tab:code003}
\end{table}

At iteration 13, JSS-CODE-003 had 429 false positives on the full
corpus. Root-cause analysis identified six distinct FP patterns:

\begin{enumerate}
  \item \emph{Hyphenated compound names}: identifiers like
    \verb|L-BFGS-B| and \verb|p-value| inside
    \code{\textbackslash code\{\}} triggered the binary-operator
    check. These hyphens are part of the name, not operators.
  \item \emph{Argument-name-only patterns}: documentation snippets
    like \code{\textbackslash code\{basis=\}} (a function argument
    name with no value, used in documentation) triggered the
    assignment-operator check.
  \item \emph{Array indexing commas}: expressions like
    \code{\textbackslash code\{a[2:5,2:5]\}} triggered the comma
    check. Array-indexing commas do not require surrounding spaces in
    \proglang{R} style.
  \item \emph{R comment lines in verbatim}: lines beginning with
    \verb|##| inside Sinput blocks were scanned for operators. A
    comment like \verb|## Pre-processing| contains a hyphen that is
    not a binary operator.
  \item \emph{Trailing-comma display notation}: closing lines of
    multi-line function calls, such as \code{\textbackslash
    code\{f(a, b),\}}, triggered the comma check on the trailing
    punctuation comma.
  \item \emph{Line-number drift}: the function that stripped Soutput
    blocks to preserve only Sinput lines was not preserving line
    count, causing violations to be reported at incorrect positions
    and thereby misattributed during classification.
\end{enumerate}

Each root cause received a targeted fix during iterations 14--15.
After all six fixes, JSS-CODE-003 reached 99.1\% precision.

\textbf{AI classification bias.} An important complication arose
during iteration 14. The AI classifier consistently labelled
\code{\textbackslash code\{param=value\}} inline patterns as false
positives, reasoning that documentation snippets like ``use
\code{kernel=TRUE} to enable \ldots{}'' are harmless prose references
that need not follow operator-spacing style. Human review of 191
such AI-labelled FPs found that 179 (94\%) were in fact genuine JSS
violations: JSS requires spaces around \code{=} in all function
arguments, including in inline documentation references.

This batch reclassification --- 179 TPs restored from FPs --- raised
JSS-CODE-003's measured precision from 72.9\% to 99.1\% in a single
iteration, demonstrating that AI classification can introduce a
systematic error when the classifier applies a plausible but incorrect
interpretation of the rule being evaluated. AI classification is a
useful first pass for volume but cannot substitute for human ground
truth when the rule's intent is subtle.

\subsection{AI-assisted code generation}
\label{sec:aigen}

Beyond classification, Claude Sonnet 4.6 contributed directly to
implementation throughout the project: rule logic, test code,
documentation, and the present manuscript. In practice, the most
effective workflow for each new rule was:

\begin{enumerate}
  \item Describe the rule in natural language, providing three to five
    representative examples of both genuine violations and similar-looking
    patterns that should \emph{not} fire.
  \item Ask the AI to propose an implementation strategy, discuss
    edge cases, and write the rule function.
  \item Review and extend the AI-generated tests to cover all branches,
    then run coverage to confirm 100\%.
  \item After the unit test suite passes: run the eval loop on the corpus
    to catch any patterns not covered by the synthetic tests.
\end{enumerate}

The AI was consistently effective at implementing rule logic once the
scope was clearly defined by examples. Its main limitation was
overconfidence in early implementations: initial rule functions
handled the motivating examples correctly but contained latent
false-positive patterns that only appeared on the full corpus.
Framing the task as ``propose an implementation \emph{and} describe
at least three ways it might fire incorrectly'' substantially improved
the quality of first drafts.

%%--------------------------------------------------------------------
\section{Challenges and retrospective lessons}
\label{sec:lessons}
%%--------------------------------------------------------------------

We reflect here on the development process, identifying what worked
well, where difficulties arose, and what we would do differently if
starting again. This section is intended as practical guidance for
others building similar tools.

\subsection{What worked well}
\label{sec:worked}

\textbf{AST-based analysis.} The decision to parse \LaTeX{} into an
AST rather than applying regular expressions to raw text eliminated an
entire class of false positives. Rules do not fire inside verbatim
environments, comment text, or macro arguments where the matched
pattern has a different semantic role. Retrofitting AST analysis into
a regex-based tool later in development is substantially harder than
designing for it from the start.

\textbf{100\% branch coverage as a policy.} Enforcing 100\% branch
coverage on all rule modules forced the writing of tests for uncommon
but important branches. Several bugs were caught because coverage
analysis revealed an untested branch, not because a test had been
written to exercise it. The policy also served as an objective
completion criterion for each task, reducing ambiguity about when
implementation was ``done''.

\textbf{The evaluation-improvement loop.} The iterative
scan-classify-report-fix loop was indispensable. Without it,
JSS-CODE-003 would have shipped at less than 65\% precision despite
passing all unit tests. No amount of synthetic test-case coverage
would have anticipated all six FP patterns discovered on real
manuscripts. The same pattern applies to any rule-based tool checking
structured text: synthetic tests establish correctness on known cases;
real-data evaluation discovers the cases that were not anticipated.

\textbf{Human review as ground truth.} Maintaining a
human-classified ground truth was critical for detecting AI classifier
bias. The database records the source of each classification decision,
making it possible to audit and revise earlier AI decisions when a
systematic error is discovered --- as happened with CODE-003.

\subsection{What proved more difficult}
\label{sec:harder}

\textbf{Synthetic tests are insufficient as the sole validation
method.} Unit tests written against a handful of synthetic \LaTeX{}
snippets do not reflect the diversity of real manuscripts.
JSS-CODE-003 passed all its unit tests but had less than 65\%
precision on real papers. Had we introduced even a small sample of
real manuscript snippets into the initial test suite --- say, twenty
genuine code blocks from existing JSS papers --- several of the six
FP patterns would have been caught during initial development rather
than during the evaluation loop.

\textbf{AI classification bias requires vigilance.} The AI classifier
introduced a systematic error for JSS-CODE-003 by applying a
plausible but incorrect interpretation of the rule. When AI-classified
precision looks unexpectedly low or high for a rule, it is worth
manually reviewing a random sample of 10--20 cases before accepting
the reported figure. A disagreement rate of even 5\% between AI and
human classifiers can distort precision estimates enough to mask or
manufacture a problem.

\textbf{Initial rule scope was too broad.} Several rules ---
JSS-CODE-003 most prominently --- were initially written to match
the motivating pattern broadly, on the assumption that edge cases
would be rare. In practice, real manuscripts exhibit far more
diversity than synthetic examples suggest, and broad rules require
multiple tightening iterations. Writing rules more narrowly from the
start, informed by analysis of real examples, reduces the number of
subsequent fix passes. A useful heuristic: if you cannot describe
in one sentence exactly what the rule does \emph{not} fire on, it is
probably too broad.

\textbf{Preprocessing bugs and line-number drift.} Preprocessing steps
such as stripping Soutput blocks from Sweave files introduce
line-number drift when stripped lines are not replaced with blank
lines of matching count. This class of bug is difficult to detect
from unit tests because synthetic tests rarely mix
content-stripping with violation line-number reporting. The root
cause took two iterations to diagnose: the first fix corrected the
stripping logic, but did not restore the correct line count;
the second fix preserved blank-line count explicitly.

\textbf{\pkg{pylatexenc} idiosyncrasies.} The library's AST node
positions use 0-based column numbers (unlike the 1-based columns
expected by most editors), and some macros with non-standard argument
specifications require special handling. These quirks caused
off-by-one bugs in early versions that were caught during integration
testing but required careful reading of the library documentation.

\textbf{Specification overhead for small changes.} The
specification-driven workflow was valuable for the five major feature
increments but felt heavy for the many small targeted fixes that arose
during the evaluation loop. Maintaining up-to-date spec/plan/tasks
documents for a five-line rule fix added overhead with little
benefit. A lighter-weight process for incremental fixes --- such as a
short comment in the commit message referencing the root cause ---
would have been more proportionate.

\subsection{Recommendations for similar projects}
\label{sec:recommendations}

\begin{enumerate}
  \item \emph{Build the evaluation loop before the first rule.} The
    scan-classify-report infrastructure is cheap to build relative to
    the rules themselves, and having it from the start means every
    rule is validated against real data immediately.
  \item \emph{Collect real examples before writing rules.} For each
    planned rule, find 10--20 real-manuscript examples of genuine
    violations and similar-looking non-violations before writing any
    code. These examples bound the rule's scope and seed the test
    suite with realistic cases.
  \item \emph{Treat AI classification as a first pass, not ground
    truth.} AI classifiers are useful for handling volume but should
    be spot-checked on every rule where the rule's intent is
    semantically subtle.
  \item \emph{Use AST analysis if the input format has a parser.}
    Regex matching on raw text is fast to implement but produces
    fragile rules. If a mature parser exists for the input format,
    invest the upfront cost of learning its API.
  \item \emph{Write narrow rules, then widen.} Starting narrow gives
    low false-positive rates from the start. Expanding a rule as the
    corpus reveals underspecified cases is safer than starting broad
    and tightening.
  \item \emph{Pin corpus content for reproducibility.} A corpus that
    changes between evaluation runs makes it impossible to distinguish
    genuine precision improvements from corpus drift. Record exact
    package versions and file hashes from the start.
\end{enumerate}

%%--------------------------------------------------------------------
\section{Implemented rules}
\label{sec:rules}
%%--------------------------------------------------------------------

\pkg{jss-style-checker} implements 53 rules in 16 categories, of
which three rules (\code{JSS-RMD-PREAMBLE-001},
\code{JSS-RMD-CAP-002}, \code{JSS-RMD-CITE-002}) apply exclusively
to \code{.Rmd} files. Table~\ref{tab:rules} gives a complete summary.

\begingroup\small
\begin{longtable}{llcp{6.8cm}}
  \caption{Complete list of rules implemented in
    \pkg{jss-style-checker}~0.3.0.
    Severity: E~=~error, W~=~warning.
    Rules prefixed \code{JSS-RMD-} apply exclusively to
    \code{.Rmd} files; rules without that prefix apply to
    \code{.tex}/\code{.Rnw} files (or to all formats for BibTeX
    rules). Rule numbers in some categories start above~001 because
    lower numbers were reserved for deprecated or planned rules.}
  \label{tab:rules} \\
  \hline
  Rule ID & Category & Sev. & Description \\
  \hline
  \endfirsthead
  \multicolumn{4}{l}{\small\textit{(Table~\ref{tab:rules} continued)}} \\
  \hline
  Rule ID & Category & Sev. & Description \\
  \hline
  \endhead
  \hline
  \endfoot
  \hline
  \endlastfoot
  \multicolumn{4}{l}{\textit{Preamble macros}} \\
  JSS-PREAMBLE-001 & Preamble & E & Missing \code{\textbackslash Plainauthor\{\}} \\
  JSS-PREAMBLE-002 & Preamble & E & Missing \code{\textbackslash Plaintitle\{\}} \\
  JSS-PREAMBLE-003 & Preamble & E & Missing \code{\textbackslash Abstract\{\}} \\
  JSS-PREAMBLE-004 & Preamble & E & Missing \code{\textbackslash Keywords\{\}} \\
  JSS-PREAMBLE-005 & Preamble & E & Missing \code{\textbackslash Plainkeywords\{\}} \\
  JSS-PREAMBLE-006 & Preamble & E & Missing \code{\textbackslash Address\{\}} \\
  JSS-RMD-PREAMBLE-001 & Preamble & E & \code{.Rmd} YAML frontmatter missing required field \\
  \multicolumn{4}{l}{\textit{Preamble format}} \\
  JSS-PRE-007 & Preamble Fmt & W & \code{\textbackslash Plain*} macro contains \LaTeX{} markup \\
  JSS-PRE-008 & Preamble Fmt & W & Keyword capitalized (non-first keyword) \\
  JSS-PRE-009 & Preamble Fmt & W & \code{\textbackslash Keywords\{\}} is empty \\
  JSS-PRE-011 & Preamble Fmt & W & \code{\textbackslash author\{\}} uses \code{\textbackslash and} instead of \code{\textbackslash And} \\
  JSS-PRE-012 & Preamble Fmt & W & \proglang{R}/Sweave manuscripts missing Sweave preamble comment \\
  \multicolumn{4}{l}{\textit{Semantic markup and math operators}} \\
  JSS-MARKUP-001 & Markup & W & Bare programming language name \\
  JSS-MARKUP-002 & Markup & W & Bare software package name \\
  JSS-MARKUP-010 & Markup & W & Bare statistical operator in math mode \\
  \multicolumn{4}{l}{\textit{Capitalisation}} \\
  JSS-CAP-002 & Capitalisation & W & Title case in section heading \\
  JSS-CAP-003 & Capitalisation & W & BibTeX title field not in title case \\
  JSS-RMD-CAP-002 & Capitalisation & W & Title case in \code{.Rmd} Markdown heading \\
  \multicolumn{4}{l}{\textit{Citations}} \\
  JSS-CITE-001 & Citations & E & Bare \code{\textbackslash cite\{\}} in parenthetical context \\
  JSS-CITE-002 & Citations & W & Bare author-year pattern in prose \\
  JSS-CITE-003 & Citations & W & \code{\textbackslash pkg\{\}} has no matching bib entry \\
  JSS-RMD-CITE-002 & Citations & W & Bare author-year pattern in \code{.Rmd} prose \\
  \multicolumn{4}{l}{\textit{Document structure}} \\
  JSS-STRUCT-001 & Structure & W & Float \code{\textbackslash label\{\}} never referenced \\
  JSS-STRUCT-002 & Structure & W & \code{\textbackslash caption\{\}} appears before float content \\
  JSS-STRUCT-003 & Structure & W & Formatting command inside \code{\textbackslash caption\{\}} \\
  JSS-STRUCT-004 & Structure & W & \code{\textbackslash footnote\{\}} inside tabular \\
  \multicolumn{4}{l}{\textit{Cross-references}} \\
  JSS-XREF-001 & Cross-refs & W & Figure missing \code{\textbackslash label\{\}} \\
  JSS-XREF-002 & Cross-refs & W & Table missing \code{\textbackslash label\{\}} \\
  JSS-XREF-003 & Cross-refs & W & Equation missing \code{\textbackslash label\{\}} \\
  \multicolumn{4}{l}{\textit{Code line width}} \\
  JSS-WIDTH-001 & Code width & W & Code line exceeds column-width threshold (default 70) \\
  \multicolumn{4}{l}{\textit{Code style}} \\
  JSS-CODE-001 & Code style & W & Non-standard indentation in verbatim block \\
  JSS-CODE-002 & Code style & W & Comment-only line in verbatim block \\
  JSS-CODE-003 & Code style & W & Missing spaces around operators in code \\
  \multicolumn{4}{l}{\textit{BibTeX completeness}} \\
  JSS-BIB-001 & BibTeX & E & \code{@article} missing required fields \\
  JSS-BIB-002 & BibTeX & E & \code{@book} missing required fields \\
  JSS-BIB-003 & BibTeX & E & \code{@inproceedings}/\code{@incollection} missing fields \\
  \multicolumn{4}{l}{\textit{Bibliography references}} \\
  JSS-REF-001 & References & W & First word after colon in title not capitalized \\
  JSS-REF-002 & References & W & Journal or proceedings name abbreviated \\
  JSS-REF-003 & References & W & Bare language/package name in BibTeX text field \\
  JSS-REF-004 & References & W & Entry with 6+ authors lacks \code{\textbackslash shortcites} \\
  JSS-REF-005 & References & W & \proglang{R} package entry missing CRAN URL \\
  JSS-REF-006 & References & W & Edition field not in ordinal format (2nd, 3rd, \ldots) \\
  \multicolumn{4}{l}{\textit{Naming conventions}} \\
  JSS-NAME-001 & Naming & W & Incorrect software name casing in text \\
  JSS-NAME-002 & Naming & W & Incorrect journal name in bib \\
  JSS-NAME-003 & Naming & W & Incorrect publisher name in bib \\
  \multicolumn{4}{l}{\textit{Typography}} \\
  JSS-TYPO-001 & Typography & W & \verb|e.g.| missing trailing comma \\
  JSS-TYPO-002 & Typography & W & \verb|i.e.| missing trailing comma \\
  JSS-TYPO-003 & Typography & W & Matrix transpose \verb|^T| instead of \code{\textbackslash top} \\
  \multicolumn{4}{l}{\textit{Abbreviations}} \\
  JSS-ABBR-001 & Abbrev. & W & Abbreviation used without prior expansion \\
  JSS-ABBR-002 & Abbrev. & W & Abbreviation with formatting or period-separated \\
  \multicolumn{4}{l}{\textit{House style}} \\
  JSS-STYLE-001 & House style & W & Hyphenated statistical term (\verb|p-value|, \verb|t-test|) \\
  JSS-STYLE-002 & House style & W & ``Subsection'' used as cross-reference prefix \\
  JSS-STYLE-003 & House style & W & Plain equation number without \code{\textbackslash ref} \\
  JSS-RCODE-001 & House style & W & Unquoted \code{library()}/\code{data()} argument \\
\end{longtable}
\endgroup

\subsection{Preamble macros}
\label{sec:rule-preamble}

Rules \code{JSS-PREAMBLE-001} through \code{JSS-PREAMBLE-006} check
that the JSS document class's six required metadata macros appear in
the preamble (before \code{\textbackslash begin\{document\}}):
\code{\textbackslash Plainauthor\{...\}},
\code{\textbackslash Plaintitle\{...\}},
\code{\textbackslash Abstract\{...\}},
\code{\textbackslash Keywords\{...\}},
\code{\textbackslash Plainkeywords\{...\}}, and
\code{\textbackslash Address\{...\}}.
The checker walks the AST nodes that precede the
\code{\textbackslash begin\{document\}} marker and records which of
the six macro names have been encountered. A separate
\code{error}-severity violation is emitted for each missing macro,
with a suggestion giving the correct macro name and a placeholder
argument.

\subsection{Semantic markup}
\label{sec:rule-markup}

Rules \code{JSS-MARKUP-001} and \code{JSS-MARKUP-002} require that
programming language names be wrapped in
\code{\textbackslash proglang\{\}} and software package names in
\code{\textbackslash pkg\{\}}. \pkg{jss-style-checker} maintains
curated lists of 30 programming language names (e.g.,{} \proglang{R},
\proglang{Python}, \proglang{Julia}, \proglang{MATLAB}, \proglang{Java})
and 51 package names (e.g.,{} \pkg{ggplot2}, \pkg{dplyr}, \pkg{numpy},
\pkg{pandas}, \pkg{lme4}).

The checker traverses the AST while tracking context state: text
nodes inside a \code{\textbackslash proglang\{\}} or
\code{\textbackslash pkg\{\}} macro are not re-checked (to avoid
flagging the argument of the correct markup), and text inside verbatim
environments is skipped entirely. Each bare occurrence generates a
\code{warning}-severity violation with a concrete replacement
suggestion.

\subsection{Capitalisation}
\label{sec:rule-cap}

\textbf{JSS-CAP-002} checks that \LaTeX{} section headings use
sentence case rather than title case. The checker extracts the heading
text from the macro argument and tests whether more than one word
(excluding common articles and conjunctions) begins with a capital
letter.

\textbf{JSS-CAP-003} checks BibTeX entry titles. JSS requires title
case in bibliographic entries so that styles that down-case titles
produce correct output when braces protect capitals. The checker strips
curly braces from the \code{title} field and verifies that at least
one content word is capitalised.

\subsection{Citations}
\label{sec:rule-cite}

\textbf{JSS-CITE-001} enforces natbib-style citation commands and
flags every \code{\textbackslash cite\{\}} macro, which produces the
doubly-parenthesised form \texttt{((Smith 2020))} when used inside
parentheses. Each occurrence generates an \code{error}-severity
violation with a suggestion to use \code{\textbackslash citep\{\}} or
\code{\textbackslash citet\{\}}.

\textbf{JSS-CITE-002} detects bare author-year patterns in prose
(e.g.,{} ``Smith (2020)'') where the surname matches a bib entry.

\textbf{JSS-CITE-003} checks that every \code{\textbackslash pkg\{\}}
macro has a corresponding bibliography entry in the provided
\code{.bib} file.

\subsection{Cross-references}
\label{sec:rule-xref}

\code{JSS-XREF-001} through \code{JSS-XREF-003} check that figure,
table, and equation environments carry a
\code{\textbackslash label\{\}}, and \code{JSS-STRUCT-001} checks
that defined labels are actually referenced somewhere in the document.

\subsection{Code line width}
\label{sec:rule-width}

\code{JSS-WIDTH-001} checks that lines inside verbatim-like
environments (\code{verbatim}, \code{Verbatim}, \code{Code},
\code{CodeInput}, \code{CodeOutput}, \code{Sinput}, \code{Soutput},
\code{lstlisting}) fit within the column width of the typeset page.
The default threshold is 70 characters, which is conservative for the
\code{jss.cls} document class at its default font size.

\subsection{BibTeX completeness}
\label{sec:rule-bib}

\code{JSS-BIB-001} through \code{JSS-BIB-003} check that
\code{@article}, \code{@book}, \code{@inproceedings}, and
\code{@incollection} entries carry all fields required for readers to
locate the cited work. Missing fields produce \code{error}-severity
violations pointing to the start line of the offending entry.

\subsection{Code style}
\label{sec:rule-codestyle}

\textbf{JSS-CODE-002} flags comment-only lines inside verbatim
environments. JSS style requires that code examples be self-contained
executable fragments; explanatory commentary should appear in
surrounding prose rather than inline code comments.

\textbf{JSS-CODE-003} flags missing spaces around binary arithmetic
operators or after commas in code blocks. The implementation handles
several edge cases that emerge frequently in real manuscripts
(Section~\ref{sec:evalloop}).

\subsection{Bibliography references}
\label{sec:rule-ref}

Six rules enforce JSS standards for BibTeX field content:
\code{JSS-REF-001} (capitalisation after a colon in titles),
\code{JSS-REF-002} (abbreviated journal or proceedings names),
\code{JSS-REF-003} (bare language/package names in bib text fields),
\code{JSS-REF-004} (entries with six or more authors lacking
\code{\textbackslash shortcites\{\}}),
\code{JSS-REF-005} (\proglang{R} package entries missing a CRAN URL),
and \code{JSS-REF-006} (edition field not in ordinal format).

\subsection{Naming, typography, abbreviations, and house style}
\label{sec:rule-misc}

\textbf{JSS-NAME-001} through \textbf{JSS-NAME-003} detect incorrect
casing for well-known software names, incorrect journal names, and
incorrect publisher names in bib fields.

\textbf{JSS-TYPO-001} and \textbf{JSS-TYPO-002} require trailing
commas after abbreviations such as ``e.g.,'' and ``i.e.,'';
\textbf{JSS-TYPO-003} requires \code{\textbackslash top} for the
matrix transpose symbol.

\textbf{JSS-ABBR-001} flags three-or-more letter uppercase
abbreviations that appear before their full expansion;
\textbf{JSS-ABBR-002} flags period-separated abbreviations and those
wrapped in font-change commands.

\textbf{JSS-MARKUP-010} enforces use of the predefined JSS math
operator macros \verb|\E|, \verb|\VAR|, \verb|\COV|, and \verb|\P| in
math mode.

\textbf{JSS-STYLE-001} flags hyphenated statistical terms (\verb|p-value|,
\verb|t-test|); \textbf{JSS-STYLE-002} and \textbf{JSS-STYLE-003}
flag ``Subsection'' as a cross-reference prefix and plain equation
numbers; \textbf{JSS-RCODE-001} flags unquoted \code{library()} and
\code{data()} arguments in \proglang{R} code.

\subsection[Rules for .Rmd source files]{Rules for \code{.Rmd} source files}
\label{sec:rule-rmd}

\textbf{JSS-RMD-PREAMBLE-001} checks that the YAML frontmatter block
contains the four fields required by the
\code{rticles::jss\_article} output format.
\textbf{JSS-RMD-CAP-002} and \textbf{JSS-RMD-CITE-002} are the
\proglang{R}~Markdown counterparts of \code{JSS-CAP-002} and
\code{JSS-CITE-002}, applied to Markdown headings and prose blocks
respectively.

%%--------------------------------------------------------------------
\section{Empirical validation}
\label{sec:eval}
%%--------------------------------------------------------------------

Rule-based linters face a fundamental precision challenge: rules
written to detect a pattern will inevitably fire on contexts where
the pattern has an innocent explanation. A unit test suite written
against synthetic examples validates only the cases the developer
anticipated. Testing against real manuscripts is essential to discover
the full range of edge cases. This section reports the results of
systematic empirical validation on 493 real JSS-format manuscripts.

\subsection{Corpus construction}
\label{sec:corpus}

Many \proglang{R} packages on CRAN include their JSS article as a
vignette in the package distribution, making them accessible
programmatically via the CRAN API \citep{cran}. We constructed a
corpus by collecting these vignettes, covering 493 \code{.tex} files
spanning a wide range of topics, writing styles, and \LaTeX{} practices
across packages published between approximately 2005 and 2026. A
pinned manifest records the exact package name, version, and
\proglang{SHA}-256 hash for each vignette, so that precision figures
reported here are reproducible against a fixed set of files. A
companion script (\code{fetch\_corpus.py}) downloads the corpus from
CRAN using the pinned manifest.

\subsection{Evaluation methodology}
\label{sec:eval-method}

A companion command-line tool, \code{eval-jss}, automates the
scan-classify-report cycle using an \proglang{SQLite} database
\citep{sqlite}. It provides five sub-commands:

\begin{itemize}
  \item \code{eval-jss init}: creates and initialises \code{eval/eval.db}.
  \item \code{eval-jss scan}: runs \code{jss-lint} against the corpus
    and stores each violation in the database, superseding earlier
    scans of the same file.
  \item \code{eval-jss review}: uses an AI language model to classify
    stored violations as true positives (TP) or false positives (FP).
  \item \code{eval-jss human-review}: presents uncertain or
    high-stakes violations for manual human classification.
  \item \code{eval-jss report}: computes per-rule precision, displays
    a summary table, and appends a timestamped row to a CSV history
    file.
\end{itemize}

\textbf{Precision metric.} Precision is defined as
$\mathrm{TP} / (\mathrm{TP} + \mathrm{FP})$, where TP counts
violations confirmed as genuine by review and FP counts violations
confirmed as incorrect. Only the \emph{latest scan per file} is
counted, so that stale violations from before a rule fix do not inflate
FP counts. Rules with fewer than ten reviewed violations are excluded
from threshold evaluation since small samples produce unreliable
estimates.

\textbf{Threshold.} We require at least 90\% precision for each rule
before considering it publication-ready. This threshold reflects the
judgment that a rule firing with more than one false positive in ten
will frustrate authors rather than help them.

\textbf{Classification procedure.} Violations were first submitted to
a locally-hosted AI classifier (a 30B-parameter language model). Cases
the AI marked as uncertain, and rules where AI reliability was known
to be lower, were routed to human review. All classification decisions
are stored with their source (AI or human) to support auditing and
revision.

\subsection{Reproducing the validation results}
\label{sec:reproducibility}

The precision figures in Table~\ref{tab:precision} are fully
reproducible from the supplementary material. The following steps
regenerate Table~\ref{tab:precision} from scratch:

\begin{enumerate}
  \item \textbf{Download the corpus.} Run
    \code{python fetch\_corpus.py} with the pinned manifest
    (\code{corpus\_manifest.csv}, included in the supplementary
    material). The script downloads each vignette from CRAN and
    verifies its \proglang{SHA}-256 hash against the manifest,
    aborting if any file does not match. The result is a local
    directory of 493 \code{.tex} files identical to those used in
    our evaluation.
  \item \textbf{Initialise the database.} Run
    \code{eval-jss init} to create a fresh \code{eval/eval.db}.
  \item \textbf{Scan the corpus.} Run \code{eval-jss scan} to invoke
    \code{jss-lint} on every file and populate the database with
    violations.
  \item \textbf{Load pre-classified decisions.} The supplementary
    material includes the final \code{eval.db} with all 6{,}739
    human- and AI-reviewed classification decisions. Importing this
    database skips the classification step and provides the same
    ground truth used for Table~\ref{tab:precision}. Alternatively,
    run \code{eval-jss review} followed by \code{eval-jss
    human-review} to re-classify violations independently.
  \item \textbf{Generate the report.} Run \code{eval-jss report} to
    reproduce Table~\ref{tab:precision}. The command also appends a
    timestamped row to \code{eval/report.csv}, providing an audit
    trail.
\end{enumerate}

Steps 1--3 and 5 are fully automated and deterministic. Step~4
(classification) requires either the supplied \code{eval.db} or
independent human review; AI-only reclassification may produce
slightly different figures due to the bias described in
Section~\ref{sec:evalloop}.

\subsection{Precision results}
\label{sec:precision}

After sixteen iterations of the scan-classify-report-fix cycle
(Section~\ref{sec:evalloop}), all thirty rules with at least ten
reviewed violations achieve at least 90\% precision.
Table~\ref{tab:precision} gives the complete results, sorted by
precision.

\begin{table}[ht]
  \centering
  \small
  \begin{tabular}{lrrr}
    \hline
    Rule & TP & FP & Precision \\
    \hline
    JSS-BIB-002      &    9 &  1 & 90.0\% \\
    JSS-REF-006      &   28 &  3 & 90.3\% \\
    JSS-RMD-CAP-002  &   11 &  1 & 91.7\% \\
    JSS-TYPO-003     &  138 & 11 & 92.6\% \\
    JSS-XREF-002     &   14 &  1 & 93.3\% \\
    JSS-MARKUP-010   &  137 &  9 & 93.8\% \\
    JSS-CITE-003     &  316 & 18 & 94.6\% \\
    JSS-RCODE-001    &   18 &  1 & 94.7\% \\
    JSS-CAP-002      &  324 & 16 & 95.3\% \\
    JSS-CITE-001     &  370 & 15 & 96.1\% \\
    JSS-NAME-003     &  223 &  9 & 96.1\% \\
    JSS-STRUCT-002   &   27 &  1 & 96.4\% \\
    JSS-STYLE-001    &   79 &  2 & 97.5\% \\
    JSS-REF-003      &  134 &  3 & 97.8\% \\
    JSS-BIB-001      &   62 &  1 & 98.4\% \\
    JSS-TYPO-001     &  273 &  4 & 98.6\% \\
    JSS-TYPO-002     &  393 &  5 & 98.7\% \\
    JSS-XREF-003     &  956 & 10 & 99.0\% \\
    JSS-CODE-003     &  809 &  7 & 99.1\% \\
    JSS-REF-001      &  180 &  1 & 99.4\% \\
    JSS-WIDTH-001    & 1369 &  4 & 99.7\% \\
    JSS-STRUCT-001   &  385 &  1 & 99.7\% \\
    JSS-BIB-003      &   65 &  0 & 100.0\% \\
    JSS-CODE-001     &   12 &  0 & 100.0\% \\
    JSS-NAME-001     &   14 &  0 & 100.0\% \\
    JSS-REF-002      &   38 &  0 & 100.0\% \\
    JSS-REF-004      &   97 &  0 & 100.0\% \\
    JSS-REF-005      &   25 &  0 & 100.0\% \\
    JSS-STRUCT-004   &   35 &  0 & 100.0\% \\
    JSS-XREF-001     &   74 &  0 & 100.0\% \\
    \hline
    \textbf{Total}   & 6615 & 124 & \textbf{98.2\%} \\
    \hline
  \end{tabular}
  \caption{Per-rule precision on the 493-manuscript corpus
    (batch-037, 2026-04-21). Only rules with at least ten reviewed
    violations are included. Sorted by precision ascending.}
  \label{tab:precision}
\end{table}

The corpus contains 6{,}615 true-positive violations and 124 false
positives, for an overall precision of 98.2\%. Eight rules achieve
100\% precision (zero false positives); the median precision across
all thirty rules is 99.0\%.

\subsection{Discussion of residual false positives}
\label{sec:eval-discussion}

The three rules with the lowest precision (all above 90\%) share a
common structure: they detect patterns that are \emph{mostly} JSS
violations but also arise legitimately in specialised document types
or non-English field content.

\textbf{JSS-BIB-002} (90.0\%, 1 FP) flags \code{@book} entries
missing required fields. The single false positive is a known
\pkg{bibtexparser} parser limitation: a malformed entry-type
declaration in one corpus file causes the parser to misclassify
a non-\code{@book} entry as \code{@book}, after which the rule
fires incorrectly. The rule logic itself is correct; the FP
reflects a parser edge case rather than a rule error.

\textbf{JSS-REF-006} (90.3\%, 3 FP) flags \code{edition} fields not
in ordinal format. The three false positives are foreign-language
editions (e.g.,{} German ordinals such as ``3.\ Auflage'') that are
correct in their original language but do not match the expected
English pattern.

\textbf{JSS-RMD-CAP-002} (91.7\%, 1 FP) flags \code{.Rmd} headings
not in sentence case. The false positive is a heading containing a
multi-word proper noun that resembles title case but is correctly
capitalised.

These residual false positives reflect fundamental limits of
pattern-based analysis: the correct decision in such edge cases
requires context that cannot be recovered from the document source
alone.

\subsection{Rules below the classification threshold}
\label{sec:eval-lowsample}

Twenty-three of the 53 rules have fewer than ten reviewed violations
in the corpus and therefore receive no precision estimate in
Table~\ref{tab:precision}. This reflects corpus composition rather
than rule failure: because all 493 papers are \proglang{R} package
vignettes, rules that target uncommon structural patterns or
non-\proglang{R} contexts fire infrequently. For example,
\code{JSS-STRUCT-003} (formatting commands inside captions) and
\code{JSS-PRE-009} (empty \code{\textbackslash Keywords}) fired on
fewer than ten papers each; \code{JSS-STYLE-003} (plain equation
numbers) and \code{JSS-CITE-002} (bare author-year in prose)
occurred but had too few classified instances for a reliable
estimate. As the corpus grows to cover more paper types these rules
will accumulate enough reviewed violations to enter the threshold
evaluation.

\subsection{Scope of the evaluation}
\label{sec:eval-scope}

\textbf{Corpus composition.} All 493 papers were drawn from CRAN
package vignettes and therefore describe \proglang{R} software. The
precision estimates in Table~\ref{tab:precision} transfer directly to
language-agnostic rules (e.g.,{} \code{JSS-CAP-002},
\code{JSS-CITE-001}, \code{JSS-TYPO-001}, \code{JSS-WIDTH-001}).
For rules that are \proglang{R}-specific by design ---
\code{JSS-CODE-003} (operator spacing in \proglang{R} code),
\code{JSS-RCODE-001} (unquoted \code{library()} arguments),
\code{JSS-PRE-012} (Sweave preamble comment) --- the estimates
characterise exactly the population of manuscripts to which the rule
applies. Authors submitting manuscripts about \proglang{Python},
\proglang{Julia}, or other non-\proglang{R} tools should expect
similar precision on the language-agnostic rules; the
\proglang{R}-specific rules will not fire at all.

\textbf{Precision vs.\ recall.} The evaluation measures
\emph{precision} (what fraction of reported violations are genuine)
but not \emph{recall} (what fraction of all genuine violations are
reported). Obtaining ground-truth recall would require a human expert
to read all 493 papers and exhaustively annotate every true
violation --- a prohibitively expensive process. Precision is the
appropriate primary metric for a tool whose main cost is false alarms
that waste authors' time. A rule with low recall is merely
incomplete; a rule with low precision actively misleads. The tool's
design goal is therefore high precision at the cost of potentially
missing some violations, not comprehensive detection at the cost of
noise. Users should be aware that a clean \code{jss-lint} run does
not certify full style compliance; it certifies that the patterns the
tool checks are not present.

%%--------------------------------------------------------------------
\section{Usage}
\label{sec:usage}
%%--------------------------------------------------------------------

\subsection{Installation}
\label{sec:install}

\pkg{jss-style-checker} requires \proglang{Python}~3.10 or later
and is installable from PyPI:

\begin{CodeInput}
pip install jss-style-checker
\end{CodeInput}

For development from source:

\begin{CodeInput}
git clone https://github.com/author/jss-style-checker
cd jss-style-checker
python3 -m venv .venv
source .venv/bin/activate
pip install -e ".[dev]"
\end{CodeInput}

\subsection{Basic usage}
\label{sec:basic-usage}

The primary entry point is the \code{jss-lint} command, which accepts
\code{.tex}, \code{.Rnw}, \code{.Rmd}, and \code{.bib} files:

\begin{CodeInput}
jss-lint paper.tex refs.bib        # LaTeX manuscript
jss-lint paper.Rnw refs.bib        # Sweave/Knitr manuscript
jss-lint paper.Rmd refs.bib        # R Markdown manuscript
\end{CodeInput}

If no violations are found the tool exits with code~0 and produces no
output. When violations are found the tool exits with code~1 and
prints a structured report. Exit code~2 indicates an invocation error
or a fatal parse failure.

\subsection{Author mode output}
\label{sec:author-output}

In the default author mode (terminal format), \code{jss-lint}
produces an annotated violation list:

\begin{verbatim}
paper.tex
  12:1  error    [JSS-PREAMBLE-001]
        Missing \Plainauthor{} macro
     -> Add \Plainauthor{name} before \begin{document}
  45:1  warning  [JSS-MARKUP-001]
        Bare language name "Python"
     -> Replace with \proglang{Python}
  78:5  error    [JSS-CITE-001]  Bare \cite{} command
     -> Use \citep{key} or \citet{key}

3 violations (2 errors, 1 warning)
\end{verbatim}

\subsection{Reviewer mode output}
\label{sec:reviewer-output}

Reviewer mode produces a per-category compliance table:

\begin{CodeOutput}
Compliance: Journal of Statistical Software
paper.tex / refs.bib

Category              Status    Violations
-------------------   ------    ----------
Preamble macros       FAIL           1
Preamble format       PASS           0
Semantic markup       FAIL           1
Math operators        PASS           0
Capitalisation        PASS           0
Citations             FAIL           1
Document structure    PASS           0
Cross-references      PASS           0
Code line width       PASS           0
Code style            PASS           0
BibTeX completeness   PASS           0
References            PASS           0
Naming conventions    PASS           0
Typography            PASS           0
Abbreviations         PASS           0
House style           PASS           0

Overall: 81.3% compliant (13 / 16 categories)
\end{CodeOutput}

\subsection{JSON output}
\label{sec:json-output}

Machine-readable output is produced with \code{--output json}.
The JSON schema (version 1.0) contains a \code{violations} array,
a \code{categories} summary, and a \code{compliance\_percentage} field.
This output is suitable for consumption by CI/CD systems, custom
dashboards, or programmatic post-processing in any language.

\subsection{Configuration file}
\label{sec:config-file}

Projects that consistently use the same options can place a
\code{.jss-lint.toml} file in the working directory:

\begin{Code}
journal = "jss"
mode = "reviewer"
output = "html"
ignore-rules = []
\end{Code}

CLI flags override the configuration file, which in turn overrides
built-in defaults.

\subsection{Auto-fix mode}
\label{sec:fix}

Many violations have a single canonical correction that can be applied
mechanically. The \code{--fix} flag instructs \code{jss-lint} to
apply these corrections directly to the source file:

\begin{CodeInput}
jss-lint --fix paper.tex refs.bib
\end{CodeInput}

Files are rewritten atomically (via \code{tempfile} and
\code{os.replace}) so that a crash during writing cannot leave a
partial file on disk. To preview what would change without touching
any files, combine \code{--fix} with \code{--dry-run}:

\begin{CodeInput}
jss-lint --fix --dry-run paper.tex refs.bib
\end{CodeInput}

Rules that currently support auto-fix are \code{JSS-MARKUP-001} (bare
programming language names $\to$ \code{\textbackslash proglang\{...\}})
and \code{JSS-NAME-001} (incorrect software name casing). Further
fixable rules are planned for future releases.

\subsection{API reference}
\label{sec:api}

\pkg{jss-style-checker} exposes a stable public API for integration
into \proglang{Python} CI workflows or editorial tools:

\begin{Code}
from texlint.core.parser import (
    parse_tex_file, parse_bib_file,
    parse_rnw_file, parse_rmd_file,
)
from texlint.api import ParsedDocument, ToolConfig
from texlint.core.engine import run

tex = parse_tex_file("paper.tex")
bib = parse_bib_file("refs.bib")
doc = ParsedDocument(tex_files=[tex], bib_files=[bib])
config = ToolConfig(
    journal="jss",
    mode="author",
    output_format="terminal",
    ignore_rules=frozenset(),
)
report = run(doc, config)
print(f"{report.compliance_percentage:.0f}% compliant")
for v in report.violations:
    print(f"  {v.line}:{v.column}  {v.rule_id}")
\end{Code}

%%--------------------------------------------------------------------
\section{Extensibility: adding journal plugins}
\label{sec:extensibility}
%%--------------------------------------------------------------------

\pkg{jss-style-checker} discovers journal rule modules at runtime via
\proglang{Python} entry points \citep{importlib_metadata}, using the
group name \code{texlint.journals}. Adding a new journal requires no
changes to the core package.

\subsection{Creating a plugin}
\label{sec:plugin-create}

A minimal journal plugin consists of a \proglang{Python} package that
subclasses \code{JournalRuleModule} and declares an entry point:

\begin{verbatim}
from texlint.api import (
    JournalRuleModule, Rule, RuleCategory, Violation
)

def _check_abstract_length(doc):
    for tex in doc.tex_files:
        pass  # implementation here

class MyJournalModule(JournalRuleModule):
    journal_id = "myjournal"
    display_name = "My Statistics Journal"
    version = "0.1.0"

    def rules(self):
        return [
            Rule(
                id="MJ-ABSTRACT-001",
                description="Abstract exceeds 150 words",
                check=_check_abstract_length,
            )
        ]

    def categories(self):
        return [
            RuleCategory(
                id="abstract",
                display_name="Abstract",
                rules=self.rules(),
            )
        ]
\end{verbatim}

\subsection{Registering the plugin}
\label{sec:plugin-register}

The plugin is registered in \code{pyproject.toml} under the
\code{texlint.journals} entry-point group:

\begin{verbatim}
[project.entry-points."texlint.journals"]
myjournal = "texlint_myjournal:MyJournalModule"
\end{verbatim}

After \code{pip install -e .}, the plugin is immediately available:

\begin{CodeInput}
jss-lint --journal myjournal paper.tex
\end{CodeInput}

The plugin system is tested in the \pkg{jss-style-checker} test suite
using a stub journal module that exercises the full discovery and
execution path.

%%--------------------------------------------------------------------
\section{Comparison with existing tools}
\label{sec:comparison}
%%--------------------------------------------------------------------

Table~\ref{tab:comparison} compares \pkg{jss-style-checker} with
the closest related tools, including both dedicated \LaTeX{} checkers
and online academic writing assistants.

\begin{table}[ht]
  \centering
  \begin{tabular}{lllll}
    \hline
    Tool & Type & JSS & AST & Offline \\
    \hline
    \pkg{chktex} \citep{chktex} & CLI & No & No & Yes \\
    \pkg{lacheck} \citep{lacheck} & CLI & No & Partial & Yes \\
    \pkg{textlint} \citep{texlint_js} & CLI & No & Yes & Yes \\
    \pkg{SciSpace} \citep{scispace} & Online & No & No & No \\
    \pkg{Writefull} \citep{writefull} & Online & No & No & No \\
    \pkg{jss-style-checker} & CLI & Yes & Yes & Yes \\
    \hline
  \end{tabular}
  \caption{Comparison of style-checking tools relevant to JSS authors.
    JSS~=~implements JSS-specific rules; AST~=~AST-based analysis;
    Offline~=~runs without internet access.}
  \label{tab:comparison}
\end{table}

\pkg{chktex}~\citep{chktex} and \pkg{lacheck}~\citep{lacheck} check
general \LaTeX{} style but know nothing about JSS conventions. Neither
can check that programming language names use
\code{\textbackslash proglang\{\}}, that preamble macros are present,
or that natbib citation commands are used correctly.

\pkg{textlint}~\citep{texlint_js} is a mature, plugin-based linting
framework for text and markup in the \proglang{JavaScript} ecosystem.
It uses an AST-based approach and supports a rich plugin system.
However, it requires \proglang{Node.js} and, to our knowledge, no
JSS-specific plugin has been published.

Online tools such as \pkg{SciSpace}~\citep{scispace} and
\pkg{Writefull}~\citep{writefull} provide AI-powered writing
assistance including grammar checking and general formatting guidance.
These tools are valuable for prose quality but do not implement
journal-specific \LaTeX{} validation rules. They also require
uploading manuscript content to third-party servers, which may be
unsuitable for unpublished work, and their output is not structured
as a machine-readable list of line-numbered violations.

To our knowledge, \pkg{jss-style-checker} is the only tool that
combines JSS-specific rules with AST-based analysis, a plugin
architecture, empirically validated precision, and fully offline
operation, in a \proglang{Python} package that fits naturally into
the workflow of authors already using \proglang{Python}-based tooling.

%%--------------------------------------------------------------------
\section{Conclusions and future work}
\label{sec:conclusions}
%%--------------------------------------------------------------------

We have presented \pkg{jss-style-checker}, a deterministic, AST-based
style checker for JSS manuscripts. The tool implements 53 rules across
16 categories, supports three output formats and two audience modes,
and accepts \code{.tex}, \code{.Rnw} (Sweave/Knitr), and \code{.Rmd}
(\proglang{R}~Markdown) source files. A \code{--fix} flag applies
deterministic auto-corrections in place (with \code{--dry-run} for
preview), and a journal plugin architecture enables other journals to
reuse the core infrastructure. Empirical evaluation on 493 real JSS
manuscripts confirms that all thirty rules with sufficient reviewed
data achieve at least 90\% precision, with an overall precision of
98.2\% across 6{,}739 reviewed violation instances.

The development followed an iterative methodology combining
specification-driven AI-assisted development with a real-world
precision evaluation loop. We have documented this methodology in
detail (Section~\ref{sec:aidev}) and reflected candidly on its
strengths and limitations (Section~\ref{sec:lessons}). The key
finding for practitioners is that synthetic unit tests, however
thorough, are insufficient to validate a rule-based linter: testing
against a corpus of real documents is essential and should be
established before the first rule is written.

The self-referential nature of this submission is intentional: the
manuscript was validated with \code{jss-lint} before submission,
confirming zero violations.

Section~\ref{sec:reproducibility} gives the five-step procedure for
reproducing Table~\ref{tab:precision} from scratch. The pinned corpus
manifest, the pre-classified \code{eval.db}, and a convenience script
\code{replicate.sh} are included as supplementary material.

Several directions for future work are planned:

\begin{itemize}
  \item \textbf{Additional journals.} The plugin architecture makes it
    straightforward to add rule sets for other statistical journals,
    e.g.,{} \emph{The \proglang{R} Journal}, \emph{Biometrics}, or
    \emph{Statistics and Computing}.
  \item \textbf{Additional JSS rules.} Future versions could check
    figure caption length, label naming conventions, and
    \code{\textbackslash bibliography} placement.
  \item \textbf{Editor integrations.} Language Server Protocol (LSP)
    integration would enable real-time feedback in editors such as
    VS Code or Emacs.
  \item \textbf{Pre-commit hook.} A ready-made pre-commit hook
    configuration would allow authors to block commits that introduce
    new violations.
  \item \textbf{Expanded auto-fix coverage.} The current
    \code{--fix} implementation corrects bare language names
    (\code{JSS-MARKUP-001}) and software name casing
    (\code{JSS-NAME-001}). Extending auto-fix to further rules
    (e.g.,{} \code{\textbackslash cite} $\to$
    \code{\textbackslash citep}, typography fixes) would increase the
    proportion of violations requiring no manual editing.
  \item \textbf{Continuous corpus expansion.} The evaluation corpus
    currently covers 493 papers. Regularly rerunning the
    evaluation-improvement loop as new JSS manuscripts become
    available will continue to surface edge cases and maintain
    precision.
\end{itemize}

\section*{Acknowledgments}

The author thanks Claude Sonnet 4.6 \citep{claude2024}, a large language
model developed by Anthropic, for substantial assistance throughout the
development of \pkg{jss-style-checker}: rule logic implementation, test
authorship, documentation, and drafting of this manuscript. All technical
decisions, empirical results, and claims in this paper are the responsibility
of the human author.

\bibliography{paper}

\end{document}
